\documentclass[12pt]{article}

% IBM-style documentation packages
\usepackage[utf8]{inputenc}
\usepackage[margin=1in, includefoot, heightrounded]{geometry}
\usepackage{graphicx}
\usepackage{float}
\usepackage{longtable}
\usepackage{booktabs}
\usepackage{enumitem}
\usepackage{courier}
\usepackage{listings}
\usepackage{xcolor}
\usepackage{fancyhdr}
\usepackage{titlesec}
\usepackage{tocloft}
\usepackage{array}
\usepackage{hyperref}
\usepackage{etoolbox}

% Font setup - Use default article font, courier only for code
% \renewcommand{\familydefault}{\ttdefault}

% Header/Footer setup
\pagestyle{fancy}
\fancyhf{}
\fancyhead[L]{\small\sffamily SeaBird ISA Reference Manual}
\fancyhead[R]{\small\sffamily Version 1.1}
\fancyfoot[C]{\small\sffamily \thepage}
\renewcommand{\headrulewidth}{0.4pt}
\renewcommand{\footrulewidth}{0.4pt}

% Section formatting
\titleformat{\section}
  {\normalfont\Large\bfseries\sffamily}{\thesection}{1em}{}
\titleformat{\subsection}
  {\normalfont\large\bfseries\sffamily}{\thesubsection}{1em}{}
\titleformat{\subsubsection}
  {\normalfont\normalsize\bfseries\sffamily}{\thesubsubsection}{1em}{}

% IBM-style code listings
\lstdefinestyle{asmstyle}{
    basicstyle=\ttfamily\small,
    commentstyle=\color{gray},
    keywordstyle=\bfseries,
    numberstyle=\tiny\color{gray},
    stringstyle=\color{black},
    breaklines=true,
    showstringspaces=false,
    frame=single,
    numbers=left,
    tabsize=4,
    xleftmargin=2em,
    framexleftmargin=1.5em,
    backgroundcolor=\color{gray!10}
}

\lstset{style=asmstyle}

% Table of contents formatting
\renewcommand{\cftsecfont}{\sffamily}
\renewcommand{\cftsubsecfont}{\sffamily}
\renewcommand{\cftsecpagefont}{\sffamily}
\renewcommand{\cftsubsecpagefont}{\sffamily}

% Hyperref setup
\hypersetup{
    colorlinks=true,
    linkcolor=black,
    urlcolor=black,
    citecolor=black,
    pdfborder={0 0 0}
}

% Custom table column types for better fitting
\newcolumntype{L}[1]{>{\raggedright\arraybackslash}p{#1}}
\newcolumntype{C}[1]{>{\centering\arraybackslash}p{#1}}
\newcolumntype{R}[1]{>{\raggedleft\arraybackslash}p{#1}}

% Adjust table spacing and appearance
\setlength{\tabcolsep}{2pt}
\renewcommand{\arraystretch}{1.1}

% Make tables use smaller font
\AtBeginEnvironment{longtable}{\footnotesize}
\AtBeginEnvironment{table}{\footnotesize}

% Better text flow - prevent overfull hboxes
\sloppy
\emergencystretch=1em

\begin{document}

% Title page
\begin{titlepage}
    \centering
    \vspace*{2cm}
    
    {\Huge\bfseries SeaBird ISA\\[0.3cm]}
    {\LARGE Reference Manual\\[0.5cm]}
    {\Large Hybrid CISC/RISC Architecture\\[0.3cm]}
    {\large Technical Specification}
    
    \vspace{2cm}
    
    {\large makateotakanawa@gmail.com\\
    Research and Development at Meisei\\[0.5cm]}
    
    \vspace{1cm}
    
    {\large December 7, 2025\\ UPDATED: January 28, 2026\\[1cm]}
    
    \vspace{1cm}
    
    \begin{center}
    \textbf{Document Version:} 1.1\\
    \textbf{Classification:} Technical Documentation
    \end{center}
    
    \vfill
    
    {\small Meisei - Naoki Saito\\
    For Use with Meisei General-Purpose Processors}
    
\end{titlepage}

\newpage

\begin{center}
\section*{Document Information}
\end{center}

\begin{table}[h]
\centering
\begin{tabular}{ll}
\toprule
\textbf{Field} & \textbf{Value} \\
\midrule
Document Title & SeaBird ISA Reference Manual \\
Version & 1.1 \\
Date & December 7, 2025 | UPDATED: January 28, 2026\\
Status & 1.1 \\
Architecture & Hybrid CISC/RISC \\
Modes Supported & Clownfish (16-bit), Tetra (32-bit), Dragonet (64-bit), Droplet (128-bit)\\
Instruction Count & 160 instructions \\
Register Count & 32 general-purpose registers (R0--R31) \\
\bottomrule
\end{tabular}
\end{table}

\vspace{1cm}

\section*{Abstract}

The SeaBird ISA architecture is a modern hybrid CISC/RISC architecture intended for use with high-performance designs. This architecture blends visible CISC machine-level instructions with a more internal execution model for micro-operations similar to a conventional RISC architecture.

The architecture supports a wide array of operating modes (either 16-bit, 32-bit, or 64-bit operation), full vectorization capabilities, and remains compatible with RISC-V.

This manual constitutes a full specification of the instruction set, encodings, register architecture, memory model, and implementation techniques for use by compiler writers, operating system designers, emulator implementors, and hardware engineers.

% Table of Contents
\tableofcontents
\newpage

% Main Content
\section{Introduction}

The SeaBird ISA is a hybrid CISC/RISC architecture designed to:

\begin{itemize}
    \item Allow highly expressive instructions for compact, human-readable assembly.
    \item Break CISC instructions into simple RISC-like micro-ops internally for high performance.
    \item Support multiple operating modes:
    \begin{itemize}
        \item Clownfish (16-bit real mode)
        \item Tetra (32-bit protected mode)
        \item Dragonet (64-bit long mode)
        \item Copycat (RISC-V compatibility mode)
    \end{itemize}
    \item Support vector and floating-point operations in addition to scalar operations.
    \item Provide a flexible register naming scheme that allows both Intel-style semantic names and Rx numeric aliases.
\end{itemize}

The architecture is designed for modern CPU pipelines, allowing micro-op fusion, vectorization, and efficient memory addressing.

\section{Operating Modes and Memory Model}

\subsection{Operating Modes}

\begin{table}[h]
\centering
\begin{tabular}{@{}llll@{}}
\toprule
\textbf{Mode} & \textbf{Bit Width} & \textbf{Purpose} & \textbf{Notes} \\ \midrule
Clownfish & 16-bit & Legacy/boot real-mode & IP=16-bit, Registers=16-bit \\
Tetra & 32-bit & Protected mode & Supports 32-bit virtual memory, \\
 &  &  & standard privilege rings \\
Dragonet & 64-bit & Long mode & 64-bit registers, supports 64-bit \\
 &  &  & addressing, RIP-relative addressing \\
Copycat & 32/64-bit & Compatibility mode & Maps R0--R31 to standard RISC-V \\
 &  &  & general-purpose registers \\ 
 Droplet & 128-bit & Hypervisor mode & Virtual wide precision mode \\ & & & for complex computing tasks\\ \bottomrule
\end{tabular}
\end{table}

\textbf{Vector Reset:} 0xCAFE (entry point for vector tables)

\subsection{Memory Model}

\begin{itemize}
    \item Flat linear memory in Tetra/Dragonet modes
    \item Memory protection and masking handled by the memory controller / MMU
    \item Instructions are independent of mode; access faults occur via MMU if privileges are violated
\end{itemize}

\subsubsection{Mode-Independent Instruction Set}

\textbf{Critical Design Principle:} All instructions in the SeaBird ISA are mode-independent. The instruction encoding, mnemonics, and behavior remain identical across all privilege levels and operating modes.

\textbf{Examples:}

\begin{lstlisting}
LD  R1, [R2]              ; loads from memory (all modes)
ST  [R3], R4              ; stores to memory (all modes)
MOV R5, R6                ; register move (all modes)
ADD R7, R8, R9            ; arithmetic (all modes)
JMP label                 ; jump (all modes)
\end{lstlisting}

Nothing about the instruction encoding changes. Nothing about the mnemonic changes. The instruction simply generates a memory request or operation.

\subsubsection{Privilege Enforcement via MMU}

The Memory Management Unit (MMU) / memory controller handles all permission checks \textit{after} instruction decode:

\textbf{Process Flow:}

\begin{enumerate}
    \item Instruction decoded (e.g., LD R1, [0xFFFF0000])
    \item Memory request issued to MMU
    \item MMU checks:
    \begin{itemize}
        \item Current privilege level (user/kernel/hypervisor)
        \item Address mapping validity
        \item Access permissions (read/write/execute)
        \item Protection attributes
    \end{itemize}
    \item MMU response:
    \begin{itemize}
        \item \textbf{Success:} Return data / complete operation
        \item \textbf{Failure:} Trigger exception (access fault, page fault, privilege violation, segmentation fault)
    \end{itemize}
\end{enumerate}

\textbf{Example Scenarios:}

\begin{table}[h]
\centering
\begin{tabular}{@{}lll@{}}
\toprule
\textbf{Mode} & \textbf{Instruction} & \textbf{Behavior} \\ \midrule
User & LD R1, [0xFFFF0000] & \textbf{Fault} if kernel-only address \\
User & ST [user\_data], R2 & \textbf{Success} if mapped \& writable \\
Kernel & LD R1, [0xFFFF0000] & \textbf{Success} (kernel memory accessible) \\
Kernel & ST [io\_port], R3 & \textbf{Success} (I/O access allowed) \\
Hypervisor & LD R1, [guest\_mem] & \textbf{Success} (VM memory accessible) \\ \bottomrule
\end{tabular}
\end{table}

\subsubsection{Mode Privilege Summary}

\textbf{User Mode:}
\begin{itemize}
    \item Can execute all instructions (LD, ST, ADD, JMP, etc.)
    \item \textit{Cannot} access kernel memory, I/O ranges, page tables, or privileged control registers
    \item Memory violations trigger exceptions handled by kernel
\end{itemize}

\textbf{Kernel Mode:}
\begin{itemize}
    \item Same instruction set as user mode
    \item \textit{Can} access privileged memory regions, I/O ports, and control registers
    \item Full access to system resources
\end{itemize}

\textbf{Hypervisor Mode (optional):}
\begin{itemize}
    \item Same instruction set as user/kernel modes
    \item \textit{Can} access VM guest memory mappings and virtualization control structures
    \item Highest privilege level
\end{itemize}

\textbf{Design Advantage:} One unified instruction set across all privilege levels. Mode determines \textit{whether access succeeds}, not \textit{what instructions are available}.

\section{Registers}

The SeaBird ISA provides a comprehensive register set with dual naming conventions, supporting both modern numeric naming (R0--R31) and legacy-style semantic aliases (AX, BX, etc.). Register width varies by operating mode.

\subsection{General-Purpose Registers (GPRs)}

\subsubsection{Primary GPRs (R0--R15)}

Available in all modes (Clownfish, Tetra, Dragonet, RISC-V):

\begin{longtable}{@{}llll@{}}
\toprule
\textbf{Physical} & \textbf{SeaBird Name} & \textbf{Legacy Alias} & \textbf{Purpose Hint} \\ \midrule
\endfirsthead
\toprule
\textbf{Physical} & \textbf{SeaBird Name} & \textbf{Legacy Alias} & \textbf{Purpose Hint} \\ \midrule
\endhead
R0 & A0 & AX & Primary accumulator \\
R1 & A1 & BX & Base/index utility \\
R2 & A2 & CX & Counter register \\
R3 & A3 & DX & Data register \\
R4 & B0 & SI & Source pointer / string ops \\
R5 & B1 & DI & Destination pointer \\
R6 & B2 & BP & Base pointer \\
R7 & B3 & SP & Stack pointer \\
R8 & C0 & R8W/R8D/R8 & Extra GPR \\
R9 & C1 & R9 & Extra GPR \\
R10 & C2 & R10 & Extra GPR \\
R11 & C3 & R11 & Extra GPR \\
R12 & D0 & R12 & Extra GPR \\
R13 & D1 & R13 & Extra GPR \\
R14 & D2 & R14 & Extra GPR \\
R15 & D3 & R15 & Extra GPR \\ \bottomrule
\end{longtable}

\subsubsection{Extended General Registers (R16--R31)}

Available in Dragonet (64-bit) mode only:

\begin{longtable}{@{}lll@{}}
\toprule
\textbf{Physical} & \textbf{SeaBird Name} & \textbf{Purpose} \\ \midrule
\endfirsthead
\toprule
\textbf{Physical} & \textbf{SeaBird Name} & \textbf{Purpose} \\ \midrule
\endhead
R16 & E0 & Environment/exception scratch \\
R17 & E1 & Syscall argument 0 \\
R18 & E2 & Syscall argument 1 \\
R19 & E3 & Syscall argument 2 \\
R20 & F0 & FP scratch 0 \\
R21 & F1 & FP scratch 1 \\
R22 & F2 & FP scratch 2 \\
R23 & F3 & FP scratch 3 \\
R24 & G0 & General purpose \\
R25 & G1 & General purpose \\
R26 & G2 & General purpose \\
R27 & G3 & General purpose \\
R28 & T0 & Temporary (µop engine) \\
R29 & T1 & Temporary (µop engine) \\
R30 & T2 & Temporary (µop engine) \\
R31 & T3 & Temporary (µop engine) \\ \bottomrule
\end{longtable}

\textbf{Note:} The T-class registers (R28--R31) are ideal for micro-op expansion and addressing-mode calculations, serving as internal scratch registers for the µop execution engine.

\subsection{Register Access Rules}

\subsubsection{Dual Naming}

All registers are accessible via multiple aliases:

\begin{itemize}
    \item \textbf{R0} = \textbf{A0} = \textbf{AX} (all refer to the same physical register)
    \item \textbf{R3} = \textbf{A3} = \textbf{DX}
    \item \textbf{R5} = \textbf{B1} = \textbf{DI}
    \item \textbf{R7} = \textbf{B3} = \textbf{SP}
\end{itemize}

\subsubsection{Width Rules Across Operating Modes}

\begin{table}[h]
\centering
\begin{tabular}{@{}lll@{}}
\toprule
\textbf{Mode} & \textbf{Register Width} & \textbf{Pointer Width} \\ \midrule
Clownfish & 16-bit & 16-bit \\
Tetra & 32-bit & 32-bit \\
Dragonet & 64-bit & 64-bit \\
Copycat & Uses R0--R31 as RV x0--x31 & 32/64-bit \\ \bottomrule
\end{tabular}
\end{table}

\subsubsection{Subregister Views}

Similar to x86-64, registers support partial access:

\begin{table}[h]
\centering
\begin{tabular}{@{}ll@{}}
\toprule
\textbf{Suffix} & \textbf{Meaning} \\ \midrule
\textbf{A0B} & Lower 8 bits (byte) \\
\textbf{A0W} & Lower 16 bits (word) \\
\textbf{A0D} & Lower 32 bits (dword) \\
\textbf{A0} & Full 64 bits (qword, Dragonet only) \\ \bottomrule
\end{tabular}
\end{table}

\textbf{Examples:}

\begin{lstlisting}
MOV A0W, #0x1234      ; move to lower 16 bits
MOV A0D, #0xDEADBEEF  ; move to lower 32 bits
MOV A0, #0xCAFEBABE   ; move to full 64-bit register
\end{lstlisting}

In Clownfish mode (16-bit):
\begin{itemize}
    \item AX refers to the full 16-bit view of R0
\end{itemize}

In Tetra mode (32-bit):
\begin{itemize}
    \item A0D or AX refers to the full 32-bit view of R0
\end{itemize}

In Dragonet mode (64-bit):
\begin{itemize}
    \item A0 refers to the full 64-bit view of R0
\end{itemize}

\subsection{Special Registers}

These registers cannot be accessed as Rn and require special instructions (similar to x86 CR/DR registers):

\begin{table}[h]
\centering
\begin{tabular}{@{}ll@{}}
\toprule
\textbf{Name} & \textbf{Purpose} \\ \midrule
IP & Instruction pointer (16/32/64-bit depending on mode) \\
FLAGS & Status flags (zero, carry, overflow, sign, etc.) \\
CFAR & CISC Frontend Address Register (RIP + prefix info) \\
UCF & Micro-op control flags \\
VECCTL & Vector mode control \\
SPTR & Shadow stack pointer (optional security feature) \\ \bottomrule
\end{tabular}
\end{table}

\textbf{Access:} Special registers are manipulated via dedicated instructions such as RDCR/WRCR or implicit behavior during control flow operations.

\section{Instruction Syntax}

\subsection{Core Syntax Rules}

\begin{lstlisting}[language={[x86masm]Assembler}]
MNEMONIC  operand1, operand2, operand3
\end{lstlisting}

\textbf{General Syntax Rules:}

\begin{itemize}
    \item \textbf{Mnemonics} are case-insensitive (ADD, add, AdD are identical)
    \item \textbf{Operands} are separated by commas
    \item \textbf{Registers} begin with R:
    \begin{itemize}
        \item R0--R15 (Clownfish mode)
        \item R0--R31 (Tetra, Dragonet, RISC-V modes)
        \item Semantic names also supported (AX, BX, SI, DI, etc.)
    \end{itemize}
    \item \textbf{Immediate values:}
    \begin{itemize}
        \item Decimal: 42
        \item Hex: 0x2A, 0xFF
        \item Binary: 0b101010, 0b1010
        \item Character: 'A'
    \end{itemize}
    \item \textbf{Comments:} ; this is a comment
    \item \textbf{Labels} end with a colon: loop\_start:
    \item \textbf{Directives} begin with . (see Assembler Directives section)
    \item \textbf{Whitespace} is ignored except inside strings
\end{itemize}

\subsection{Syntax Examples}

\begin{lstlisting}
; Case-insensitive mnemonics
ADD R1, R2, R3
add r1, r2, r3        ; equivalent to above

; Immediate values
MOVI R0, #42          ; decimal
MOVI R1, #0x2A        ; hexadecimal
MOVI R2, #0b101010    ; binary
MOVI R3, #'A'         ; character literal

; Labels and jumps
loop_start:
    ADD R1, R2, R3
    DEC R4
    JNZ loop_start    ; jump if not zero

; Comments
MOV R0, R1            ; copy R1 to R0
\end{lstlisting}

\section{Addressing Modes}

The SeaBird ISA uses CISC-style addressing modes in the programmer-visible frontend, which are internally decomposed into RISC-like micro-ops by the backend execution engine. This unified addressing model is available across all operating modes.

\begin{longtable}{@{}lll@{}}
\toprule
\textbf{Mode} & \textbf{Syntax} & \textbf{Description} \\ \midrule
\endfirsthead
\toprule
\textbf{Mode} & \textbf{Syntax} & \textbf{Description} \\ \midrule
\endhead
Register Direct & R1 & Operand in register \\
Immediate & \#42, \#0xFF & Constant value (decimal, hex, binary) \\
Register Indirect & [R2] & Memory pointed by register \\
Register + Offset & [R2+16], [R7+0x20] & Base + immediate displacement \\
Base + Index (Scaled) & [R1+R3*4], [R2+R6*8] & Base + scaled index (typical array access) \\
Base + Index + Displacement & [R4+R5*2+32] & Full complex addressing \\
PC-relative & label, label(RIP) & PC-relative (Dragonet mode only) \\
Absolute Address & [0x9000] & Direct memory address (primarily Clownfish/Tetra) \\ \bottomrule
\end{longtable}

\subsection{Addressing Mode Examples}

\begin{lstlisting}
; Register Direct
MOV R1, R2                ; copy R2 to R1

; Immediate
MOVI R3, #42              ; load immediate value 42
ADDI R4, #0xFF            ; add 255 to R4

; Register Indirect
LD R5, [R6]               ; load from address in R6
ST [R7], R8               ; store R8 to address in R7

; Register + Offset
LD R1, [R2+16]            ; load from R2 + 16
ST [R4+0x20], R3          ; store to R4 + 32

; Base + Index (Scaled)
LD R0, [R1+R3*4]          ; array access: base + index*4
ST [R2+R6*8], R5          ; store with 8-byte stride

; Base + Index + Displacement
LD R7, [R4+R5*2+32]       ; complex: base + index*2 + 32

; PC-relative (Dragonet mode)
JMP forward               ; relative to current IP
CALL function(RIP)        ; RIP-relative call

; Absolute Address (Clownfish/Tetra)
LD R1, [0x9000]           ; load from fixed address
ST [0xFFFF], R2           ; store to fixed address

forward:
    RET
\end{lstlisting}

\subsection{Micro-Op Decomposition}

Complex addressing modes are transparently broken into micro-ops:

\begin{lstlisting}
; Programmer writes:
LD R1, [R2+R3*4+32]

; Internal micro-op decomposition:
uMUL  Rtmp, R3, #4        ; scale index
uADD  Rtmp, R2, Rtmp      ; add base
uADD  Rtmp, Rtmp, #32     ; add displacement
uLOAD R1, [Rtmp]          ; perform load
\end{lstlisting}

\section{Instruction Formats}

The SeaBird ISA presents compact, expressive CISC-style instruction formats to the programmer. Internally, the backend decomposes these into RISC-like micro-ops during execution.

\subsection{CISC Frontend Formats}

\begin{longtable}{@{}lll@{}}
\toprule
\textbf{Format} & \textbf{Syntax} & \textbf{Usage} \\ \midrule
\endfirsthead
\toprule
\textbf{Format} & \textbf{Syntax} & \textbf{Usage} \\ \midrule
\endhead
R-type & MNEMONIC Rdst, Rsrc1, Rsrc2 & Register-only operations \\
 &  & (arithmetic, logic, compare, bitwise) \\
I-type & MNEMONIC Rdst, Rsrc, \#imm & Immediate arithmetic / logic \\
M-type & MNEMONIC Rdst, [mem] & Memory access (load/store) \\
 & MNEMONIC [mem], Rsrc & \\
B-type & BCOND label & Branches and jumps \\
 & JMP label / JZ label & (conditional/unconditional) \\
S-type & MNEMONIC sysop & System / control instructions \\
 &  & (interrupts, mode-switch, flags) \\
V-type & VADD V0, V1, V2 & Vector/matrix operations \\
 & VMUL V3, V1, \#5 & \\
X-type & MNEMONIC extended & Extended / microcoded instructions \\
 &  & (multi-µop guaranteed expansion) \\ \bottomrule
\end{longtable}

\subsection{Detailed Format Descriptions}

\subsubsection{R-Type (Register Only)}

Register-to-register operations for arithmetic, logic, comparison, and bitwise operations.

\begin{lstlisting}
ADD   R1, R2, R3          ; R1 = R2 + R3
SUB   R4, R5, R6          ; R4 = R5 - R6
AND   R7, R8, R9          ; R7 = R8 & R9
CMP   R1, R2              ; compare R1 and R2
MAX   R3, R4              ; R3 = max(R3, R4)
\end{lstlisting}

\subsubsection{I-Type (Immediate)}

Operations involving immediate (constant) values.

\begin{lstlisting}
ADDI  R1, R2, #42         ; R1 = R2 + 42
MOVI  R3, #0xFF           ; R3 = 255
CMPI  R4, #100            ; compare R4 with 100
SUBI  R5, #8              ; R5 = R5 - 8
\end{lstlisting}

\subsubsection{M-Type (Memory Access)}

Load and store operations with various addressing modes.

\begin{lstlisting}
LOAD  R1, [R2]            ; load from [R2]
LOAD  R1, [R2+4]          ; load from [R2+4]
LOAD  R1, [R2+R3*4]       ; load from [R2+R3*4]
STORE [R4], R5            ; store R5 to [R4]
STORE [R4+R5*2+32], R6    ; complex addressing
\end{lstlisting}

\subsubsection{B-Type (Branch / Jump)}

Control flow operations, both conditional and unconditional.

\begin{lstlisting}
JMP   label               ; unconditional jump
JZ    loop_start          ; jump if zero
JNZ   exit                ; jump if not zero
JG    greater             ; jump if greater
CALL  function            ; call subroutine
RET                       ; return from subroutine
\end{lstlisting}

\subsubsection{S-Type (System / Control)}

System-level operations for interrupts, mode switching, and privileged operations.

\begin{lstlisting}
SYSCALL                   ; system call
HLT                       ; halt processor
CLI                       ; clear interrupts
STI                       ; set interrupts
RDCR  R1, CR0             ; read control register
WRCR  CR2, R3             ; write control register
\end{lstlisting}

\subsubsection{V-Type (Vector / Matrix)}

SIMD and vector operations operating on vector registers.

\begin{lstlisting}
VADD  V0, V1, V2          ; vector add
VMUL  V3, V4, V5          ; vector multiply
VSUB  V6, V7, V8          ; vector subtract
VDIV  V0, V1, V2          ; vector divide
VADD  V3, V1, #5          ; vector-scalar add
\end{lstlisting}

\subsubsection{X-Type (Extended / Microcoded)}

Special instructions guaranteed to expand into multiple micro-ops. Used for complex operations.

\begin{lstlisting}
PUSHA                     ; push all GPRs (microcoded)
POPA                      ; pop all GPRs (microcoded)
ENTER #16                 ; create stack frame
LEAVE                     ; destroy stack frame
CPYB  R1, R2, #256        ; copy 256 bytes (microcoded loop)
\end{lstlisting}

\section{Instruction Encoding}

\subsection{Binary Encoding Overview}

The SeaBird ISA uses an x86-style modular binary encoding format. All instructions follow this layout:

\begin{center}
\texttt{[Prefix][Opcode][ModR/M][SIB][Displacement/Immediate]}
\end{center}

Each instruction may consist of the following components:

\begin{enumerate}
    \item \textbf{Prefix byte} (optional, 0--1 bytes) - Mode, size overrides, vector selectors
    \item \textbf{Opcode byte(s)} (1--3 bytes) - Instruction identifier
    \item \textbf{ModR/M byte} (optional) - Register/memory operand encoding
    \item \textbf{SIB byte} (optional) - Scaled index addressing
    \item \textbf{Immediate / Displacement} (0, 1, 4, or 8 bytes) - Constants and offsets
\end{enumerate}

This encoding scheme is:
\begin{itemize}
    \item \textbf{CISC-friendly:} Compact representation for programmers
    \item \textbf{Mode-independent:} Same encoding works across 16/32/64-bit modes
    \item \textbf{Extensible:} Easy to add vector and system instructions
    \item \textbf{RISC-backend compatible:} Decomposes cleanly into micro-ops
\end{itemize}

\subsection{Prefix Byte (Optional, 1 byte)}

The prefix byte modifies instruction behavior and is only required for overrides or special operations.

\begin{table}[h]
\centering
\begin{tabular}{@{}lll@{}}
\toprule
\textbf{Bits} & \textbf{Purpose} & \textbf{Notes} \\ \midrule
7-6 & Mode & 00=Clownfish (16-bit), 01=Tetra (32-bit), \\
 &  & 10=Dragonet (64-bit), 11=Reserved \\
5 & Address size override & Switch between default mode and 16/32/64-bit addresses \\
4 & Operand size override & Switch default operand size (16/32/64-bit) \\
3-0 & Vector/SIMD selector & Lane width, matrix/vector size, etc. \\ \bottomrule
\end{tabular}
\end{table}

\textbf{Usage:} If instruction operates in default mode, prefix may be omitted. Required only for overrides, vector operations, or cross-mode compatibility.

\subsection{Opcode Bytes (1--3 bytes)}

Opcodes are organized by functional class for efficient decoding:

\begin{table}[h]
\centering
\begin{tabular}{@{}llll@{}}
\toprule
\textbf{Class} & \textbf{Opcode Range} & \textbf{Count} & \textbf{Bytes} \\ \midrule
Data movement & 0x00--0x1F & 24 instructions & 1--2 \\
Arithmetic & 0x20--0x3F & 26 instructions & 1 \\
Logic/Bit & 0x40--0x51 & 18 instructions & 1 \\
Comparison & 0x52--0x5B & 10 instructions & 1 \\
Branch/Control & 0x5C--0x71 & 20 instructions & 1--2 \\
Memory extensions & 0x72--0x7D & 12 instructions & 2 \\
System/Privileged & 0x7E--0x8C & 15 instructions & 2--3 \\
Stack/Call & 0x8D--0x96 & 10 instructions & 1--2 \\
SIMD/Vector & 0x97--0xA5 & 15 instructions & 2 \\
Floating Point & 0xA6--0xAF & 10 instructions & 2 \\ \bottomrule
\end{tabular}
\end{table}

\textbf{Encoding Strategy:}
\begin{itemize}
    \item 1 byte: Basic arithmetic, logic, comparison, stack operations
    \item 2 bytes: Memory operations, vector operations
    \item 3 bytes: System, privileged, complex CISC instructions (e.g., LEAS, PUSHA/POPA)
\end{itemize}

\subsection{ModR/M Byte (Optional)}

The ModR/M byte encodes register and memory operands, similar to x86 but simplified.

\begin{table}[h]
\centering
\begin{tabular}{@{}ll@{}}
\toprule
\textbf{Bits} & \textbf{Meaning} \\ \midrule
7-6 & mod (addressing mode) \\
5-3 & reg (register operand or opcode extension) \\
2-0 & r/m (register/memory operand) \\ \bottomrule
\end{tabular}
\end{table}

\textbf{mod Field Values:}

\begin{table}[h]
\centering
\begin{tabular}{@{}ll@{}}
\toprule
\textbf{mod} & \textbf{Addressing Mode} \\ \midrule
00 & [Rbase] - register indirect \\
01 & [Rbase + imm8] - base + 8-bit displacement \\
10 & [Rbase + imm32] - base + 32-bit displacement \\
11 & Register direct (no memory access) \\ \bottomrule
\end{tabular}
\end{table}

\textbf{Note:} This encoding is cleaner than x86 with no legacy mode complications. It fits all 16/32/64-bit modes uniformly.

\subsection{SIB Byte (Optional)}

The SIB (Scale-Index-Base) byte enables scaled index addressing for array and structure access.

\begin{table}[h]
\centering
\begin{tabular}{@{}ll@{}}
\toprule
\textbf{Bits} & \textbf{Meaning} \\ \midrule
7-6 & scale (00=1, 01=2, 10=4, 11=8) \\
5-3 & index register (R0--R7 in 3 bits, extended via prefix) \\
2-0 & base register (R0--R7 in 3 bits, extended via prefix) \\ \bottomrule
\end{tabular}
\end{table}

\textbf{Usage:} Only present when instruction uses [base + index*scale + disp] addressing.

\subsection{Displacement / Immediate (0, 1, 4, or 8 bytes)}

Variable-length field for constants and memory offsets.

\begin{table}[h]
\centering
\begin{tabular}{@{}lll@{}}
\toprule
\textbf{Size} & \textbf{Usage} & \textbf{Example} \\ \midrule
0 bytes & No immediate/displacement & ADD R1, R2, R3 \\
1 byte & Small displacement/immediate & LD R1, [R2+16] (Clownfish) \\
4 bytes & 32-bit displacement/immediate & ADDI R3, \#0x12345678 (Tetra) \\
8 bytes & 64-bit displacement/immediate & MOVI R4, \#0xCAFEBABE (Dragonet) \\ \bottomrule
\end{tabular}
\end{table}

\subsection{Addressing Modes Summary}

\begin{longtable}{@{}lll@{}}
\toprule
\textbf{Mode} & \textbf{Syntax} & \textbf{Description} \\ \midrule
\endfirsthead
\toprule
\textbf{Mode} & \textbf{Syntax} & \textbf{Description} \\ \midrule
\endhead
Register direct & Rdst & Simple register to register \\
Memory absolute & [addr] & Direct memory access \\
Register indirect & [Rbase] & Load/store through base register \\
Base + offset & [Rbase + imm] & Small/large displacement \\
Indexed & [Rbase + Rindex*scale + imm] & For arrays / structures \\
Stack & PUSH/POP & Special SP-relative addressing \\ \bottomrule
\end{longtable}

\subsection{Encoding Examples}

\subsubsection{Example 1: Simple Arithmetic}

\begin{lstlisting}
ADD R1, R2        ; 16-bit Clownfish mode
\end{lstlisting}

\textbf{Binary Encoding:}
\begin{itemize}
    \item No prefix (default mode)
    \item Opcode: 0x20
    \item ModR/M: mod=11 (register), reg=R2 (010), r/m=R1 (001)
    \item Total: 2 bytes
\end{itemize}

\subsubsection{Example 2: Complex Addressing}

\begin{lstlisting}
LEAS R4, [R5+R6*4+16]    ; 32-bit Tetra mode
\end{lstlisting}

\textbf{Binary Encoding:}
\begin{itemize}
    \item Prefix: 0x40 (Tetra mode indicator)
    \item Opcode: 0x1E (LEAS)
    \item ModR/M: mod=10 (32-bit disp), reg=R4 (100), r/m=R5 (101)
    \item SIB: scale=10 (4x), index=R6 (110), base=R5 (101)
    \item Displacement: 4 bytes (0x00000010)
    \item Total: 9 bytes
\end{itemize}

\subsubsection{Example 3: Vector Operation}

\begin{lstlisting}
VADD V0, V1, V2          ; 128-bit vector
\end{lstlisting}

\textbf{Binary Encoding:}
\begin{itemize}
    \item Prefix: 0x08 (vector selector: 128-bit)
    \item Opcode: 0x97 (VADD)
    \item Operands: encoded in following bytes (V0, V1, V2)
    \item Total: 3--4 bytes
\end{itemize}

\subsubsection{Example 4: Immediate Load}

\begin{lstlisting}
MOVI R3, #0xDEADBEEF     ; 32-bit immediate
\end{lstlisting}

\textbf{Binary Encoding:}
\begin{itemize}
    \item No prefix (default operand size)
    \item Opcode: 0x01 (MOVI)
    \item ModR/M: mod=11, reg=000 (unused), r/m=R3 (011)
    \item Immediate: 4 bytes (0xDEADBEEF)
    \item Total: 6 bytes
\end{itemize}

\subsection{Encoding Legend}

\begin{table}[h]
\centering
\begin{tabular}{@{}ll@{}}
\toprule
\textbf{Column} & \textbf{Meaning} \\ \midrule
Instr & Mnemonic (with operands) \\
Class & Instruction class (data/arithmetic/etc.) \\
Opcode & Primary opcode byte(s) \\
ModR/M & ModR/M byte usage (if applicable) \\
SIB & SIB byte usage (if applicable) \\
Imm/Disp & Immediate/displacement byte size \\
Micro-Ops & Micro-op breakdown, internal RISC-like ops \\ \bottomrule
\end{tabular}
\end{table}

\textbf{Notes:}

\begin{itemize}
    \item Rdst, Rsrc denote registers (R0--R31)
    \item [mem] denotes memory operand
    \item Prefix byte optional (mode, operand size, vector width)
\end{itemize}

\section{Instruction Encoding \& Micro-Op Mapping}

This section provides the complete encoding details and micro-op decomposition for all 160 instructions in the SeaBird ISA.

\subsection{Data Movement (24 instructions)}

\begin{longtable}{@{}L{3.5cm}L{1.5cm}L{2cm}L{4.5cm}@{}}
\toprule
\textbf{Instruction} & \textbf{Opcode} & \textbf{Encoding} & \textbf{Micro-Operations} \\ \midrule
\endfirsthead
\toprule
\textbf{Instruction} & \textbf{Opcode} & \textbf{Encoding} & \textbf{Micro-Operations} \\ \midrule
\endhead
MOV Rdst, Rsrc & 0x00 & R-type & uMOV Rdst,Rsrc \\
MOVI Rdst, \#imm & 0x01 & I-type (1/2/4/8B) & uMOV Rdst,\#imm \\
MOVZX Rdst, Rsrc & 0x02 & R-type & uMOV Rdst,Rsrc; uZEXT Rdst \\
MOVSX Rdst, Rsrc & 0x03 & R-type & uMOV Rdst,Rsrc; uSEXT Rdst \\
MOVHI Rdst, Rsrc & 0x04 & R-type & uMOVHI Rdst,Rsrc \\
MOVLO Rdst, Rsrc & 0x05 & R-type & uMOVLO Rdst,Rsrc \\
MOVSWP Rdst, Rsrc & 0x06 & R-type & uBSWAP Rdst \\
LD Rdst, [addr] & 0x10 & M-type + disp & uLOAD Rdst,[addr] \\
ST [addr], Rsrc & 0x11 & M-type + disp & uSTORE [addr],Rsrc \\
LDI Rdst, \#imm & 0x12 & I-type + disp & uLOAD Rdst,[\#addr+imm] \\
LDB Rdst, [addr] & 0x13 & M-type (1B) & uLOAD.B Rdst,[addr] \\
LDH Rdst, [addr] & 0x14 & M-type (2B) & uLOAD.H Rdst,[addr] \\
LDW Rdst, [addr] & 0x15 & M-type (4B) & uLOAD.W Rdst,[addr] \\
LDQ Rdst, [addr] & 0x16 & M-type (8B) & uLOAD.Q Rdst,[addr] \\
STB [addr], Rsrc & 0x17 & M-type (1B) & uSTORE.B [addr],Rsrc \\
STH [addr], Rsrc & 0x18 & M-type (2B) & uSTORE.H [addr],Rsrc \\
STW [addr], Rsrc & 0x19 & M-type (4B) & uSTORE.W [addr],Rsrc \\
STQ [addr], Rsrc & 0x1A & M-type (8B) & uSTORE.Q [addr],Rsrc \\
LDP Rdst1, Rdst2, [addr] & 0x1B & M-type + pair & uLOAD Rdst1,[addr]; uLOAD Rdst2,[addr+sz] \\
STP [addr], Rsrc1, Rsrc2 & 0x1C & M-type + pair & uSTORE [addr],Rsrc1; uSTORE [addr+sz],Rsrc2 \\
LEA Rdst, [addr] & 0x1D & M-type & uADDR Rdst,[addr] \\
LEAS Rdst, [base+idx*s+d] & 0x1E & M-type + SIB & uMUL tmp,idx,s; uADD tmp,base,tmp; uADD Rdst,tmp,d \\
XCHG Rdst, Rsrc & 0x1F & R-type & uXCHG Rdst,Rsrc \\ \bottomrule
\end{longtable}

\subsection{Arithmetic (26 instructions)}

\begin{longtable}{@{}L{3cm}L{1.5cm}L{1.8cm}L{4.8cm}@{}}
\toprule
\textbf{Instruction} & \textbf{Opcode} & \textbf{Encoding} & \textbf{Micro-Operations} \\ \midrule
\endfirsthead
\toprule
\textbf{Instruction} & \textbf{Opcode} & \textbf{Encoding} & \textbf{Micro-Operations} \\ \midrule
\endhead
ADD Rdst, Rsrc & 0x20 & R-type & uADD Rdst,Rsrc \\
ADDI Rdst, \#imm & 0x21 & I-type (1/2/4/8B) & uADD Rdst,\#imm \\
SUB Rdst, Rsrc & 0x22 & R-type & uSUB Rdst,Rsrc \\
SUBI Rdst, \#imm & 0x23 & I-type (1/2/4/8B) & uSUB Rdst,\#imm \\
MUL Rdst, Rsrc & 0x24 & R-type & uMUL Rdst,Rsrc \\
MULI Rdst, \#imm & 0x25 & I-type (1/2/4/8B) & uMUL Rdst,\#imm \\
DIV Rdst, Rsrc & 0x26 & R-type & uDIV Rdst,Rsrc \\
DIVI Rdst, \#imm & 0x27 & I-type (1/2/4/8B) & uDIV Rdst,\#imm \\
MOD Rdst, Rsrc & 0x28 & R-type & uMOD Rdst,Rsrc \\
MODI Rdst, \#imm & 0x29 & I-type (1/2/4/8B) & uMOD Rdst,\#imm \\
NEG Rdst & 0x2A & R-type (unary) & uNEG Rdst \\
INC Rdst & 0x2B & R-type (unary) & uINC Rdst \\
DEC Rdst & 0x2C & R-type (unary) & uDEC Rdst \\
MULH Rdst, Rsrc & 0x2D & R-type (high bits) & uMULH Rdst,Rsrc \\
UMUL Rdst, Rsrc & 0x2E & R-type (unsigned) & uUMUL Rdst,Rsrc \\
UDIV Rdst, Rsrc & 0x2F & R-type (unsigned) & uUDIV Rdst,Rsrc \\
ADDS Rdst, Rsrc & 0x30 & R-type (saturate) & uADD.SAT Rdst,Rsrc \\
ADDU Rdst, Rsrc & 0x31 & R-type (sat uns) & uADD.USAT Rdst,Rsrc \\
SUBS Rdst, Rsrc & 0x32 & R-type (saturate) & uSUB.SAT Rdst,Rsrc \\
SUBU Rdst, Rsrc & 0x33 & R-type (sat uns) & uSUB.USAT Rdst,Rsrc \\
ABS Rdst, Rsrc & 0x34 & R-type & uABS Rdst,Rsrc \\
CLZ Rdst, Rsrc & 0x35 & R-type & uCLZ Rdst,Rsrc \\
CTZ Rdst, Rsrc & 0x36 & R-type & uCTZ Rdst,Rsrc \\
POPC Rdst, Rsrc & 0x37 & R-type & uPOPC Rdst,Rsrc \\ \bottomrule
\end{longtable}

\subsection{Logic \& Bit Operations (18 instructions)}

\begin{longtable}{@{}L{3.2cm}L{1.5cm}L{1.7cm}L{4.8cm}@{}}
\toprule
\textbf{Instruction} & \textbf{Opcode} & \textbf{Encoding} & \textbf{Micro-Operations} \\ \midrule
\endfirsthead
\toprule
\textbf{Instruction} & \textbf{Opcode} & \textbf{Encoding} & \textbf{Micro-Operations} \\ \midrule
\endhead
AND Rdst, Rsrc & 0x40 & R-type & uAND Rdst,Rsrc \\
OR Rdst, Rsrc & 0x41 & R-type & uOR Rdst,Rsrc \\
XOR Rdst, Rsrc & 0x42 & R-type & uXOR Rdst,Rsrc \\
NOT Rdst & 0x43 & R-type (unary) & uNOT Rdst \\
NAND Rdst, Rsrc & 0x44 & R-type & uNAND Rdst,Rsrc \\
NOR Rdst, Rsrc & 0x45 & R-type & uNOR Rdst,Rsrc \\
XNOR Rdst, Rsrc & 0x46 & R-type & uXNOR Rdst,Rsrc \\
SHL Rdst, Rsrc & 0x47 & R-type & uSHL Rdst,Rsrc \\
SHR Rdst, Rsrc & 0x48 & R-type & uSHR Rdst,Rsrc \\
SAR Rdst, Rsrc & 0x49 & R-type & uSAR Rdst,Rsrc \\
ROL Rdst, Rsrc & 0x4A & R-type & uROL Rdst,Rsrc \\
ROR Rdst, Rsrc & 0x4B & R-type & uROR Rdst,Rsrc \\
BSET Rdst, bit & 0x4C & R-type + imm & uBSET Rdst,bit \\
BCLR Rdst, bit & 0x4D & R-type + imm & uBCLR Rdst,bit \\
BTOG Rdst, bit & 0x4E & R-type + imm & uBTOG Rdst,bit \\
BTST Rdst, bit & 0x4F & R-type + imm & uBTST Rdst,bit \\
MASK Rdst, Rsrc, mask & 0x50 & R-type + imm & uMASK Rdst,Rsrc,mask \\
EXT Rdst, Rsrc, range & 0x51 & R-type + imm & uEXT Rdst,Rsrc,range \\ \bottomrule
\end{longtable}

\subsection{Comparison (10 instructions)}

\begin{longtable}{@{}L{2.8cm}L{1.5cm}L{1.7cm}L{4.8cm}@{}}
\toprule
\textbf{Instruction} & \textbf{Opcode} & \textbf{Encoding} & \textbf{Micro-Operations} \\ \midrule
\endfirsthead
\toprule
\textbf{Instruction} & \textbf{Opcode} & \textbf{Encoding} & \textbf{Micro-Operations} \\ \midrule
\endhead
CMP Ra, Rb & 0x52 & R-type & uCMP Ra,Rb (set flags) \\
CMPI Ra, \#imm & 0x53 & I-type & uCMP Ra,\#imm (set flags) \\
CMPS Ra, Rb & 0x54 & R-type & uCMP.S Ra,Rb (signed) \\
CMPU Ra, Rb & 0x55 & R-type & uCMP.U Ra,Rb (unsigned) \\
TST Ra, Rb & 0x56 & R-type & uTST Ra,Rb (bitwise AND, set flags) \\
TSTI Ra, \#imm & 0x57 & I-type & uTST Ra,\#imm (test bits) \\
MAX Rdst, Rsrc & 0x58 & R-type & uMAX Rdst,Rsrc (max value) \\
MIN Rdst, Rsrc & 0x59 & R-type & uMIN Rdst,Rsrc (min value) \\
SLT Rdst, Rsrc & 0x5A & R-type & uSLT Rdst,Rsrc (set if less than) \\
SGT Rdst, Rsrc & 0x5B & R-type & uSGT Rdst,Rsrc (set if greater) \\ \bottomrule
\end{longtable}

\subsection{Branch / Control Flow (20 instructions)}

\begin{longtable}{@{}L{2.8cm}L{1.5cm}L{1.7cm}L{4.5cm}@{}}
\toprule
\textbf{Instruction} & \textbf{Opcode} & \textbf{Encoding} & \textbf{Micro-Operations} \\ \midrule
\endfirsthead
\toprule
\textbf{Instruction} & \textbf{Opcode} & \textbf{Encoding} & \textbf{Micro-Operations} \\ \midrule
\endhead
JMP label & 0x5C & B-type (rel) & uJMP PC+offset \\
JMPA addr & 0x5D & B-type (abs) & uJMP addr \\
CALL label & 0x5E & B-type (rel) & uPUSH ret\_addr; uJMP PC+offset \\
CALLA addr & 0x5F & B-type (abs) & uPUSH ret\_addr; uJMP addr \\
RET & 0x60 & B-type & uPOP PC \\
JE label & 0x61 & B-type (cond) & uBRANCH\_IF ZF=1 \\
JNE label & 0x62 & B-type (cond) & uBRANCH\_IF ZF=0 \\
JG label & 0x63 & B-type (cond) & uBRANCH\_IF (SF=OF) \& ZF=0 \\
JGE label & 0x64 & B-type (cond) & uBRANCH\_IF SF=OF \\
JL label & 0x65 & B-type (cond) & uBRANCH\_IF SF!=OF \\
JLE label & 0x66 & B-type (cond) & uBRANCH\_IF (SF!=OF) | ZF=1 \\
JC label & 0x67 & B-type (cond) & uBRANCH\_IF CF=1 \\
JNC label & 0x68 & B-type (cond) & uBRANCH\_IF CF=0 \\
JO label & 0x69 & B-type (cond) & uBRANCH\_IF OF=1 \\
JNO label & 0x6A & B-type (cond) & uBRANCH\_IF OF=0 \\
JS label & 0x6B & B-type (cond) & uBRANCH\_IF SF=1 \\
JNS label & 0x6C & B-type (cond) & uBRANCH\_IF SF=0 \\
JZR reg, label & 0x6D & B-type (reg) & uBRANCH\_IF reg=0 \\
JNZR reg, label & 0x6E & B-type (reg) & uBRANCH\_IF reg!=0 \\
BRR reg & 0x6F & B-type (ind) & uJMP [reg] \\
TRAP \#imm & 0x70 & S-type & uTRAP imm (exception) \\
YIELD & 0x71 & S-type & uYIELD (scheduler hint) \\ \bottomrule
\end{longtable}

\subsection{Memory / Addressing Extensions (12 instructions)}

\begin{longtable}{@{}L{3.2cm}L{1.5cm}L{1.7cm}L{4.5cm}@{}}
\toprule
\textbf{Instruction} & \textbf{Opcode} & \textbf{Encoding} & \textbf{Micro-Operations} \\ \midrule
\endfirsthead
\toprule
\textbf{Instruction} & \textbf{Opcode} & \textbf{Encoding} & \textbf{Micro-Operations} \\ \midrule
\endhead
PREFETCH [addr] & 0x72 & M-type & uPREFETCH [addr] (hint) \\
FLUSH [addr] & 0x73 & M-type & uFLUSH [addr] (cache writeback) \\
INVIC [addr] & 0x74 & M-type & uINVAL\_ICACHE [addr] \\
INVDC [addr] & 0x75 & M-type & uINVAL\_DCACHE [addr] \\
LDX Rdst, [addr+] & 0x76 & M-type & uLOAD Rdst,[addr]; uINC addr \\
STX [addr+], Rsrc & 0x77 & M-type & uSTORE [addr],Rsrc; uINC addr \\
LDN Rdst, [addr-] & 0x78 & M-type & uDEC addr; uLOAD Rdst,[addr] \\
STN [addr-], Rsrc & 0x79 & M-type & uDEC addr; uSTORE [addr],Rsrc \\
MOVM Rdst, mask & 0x7A & X-type & Multi-reg masked move (microcoded) \\
CPYB dst, src, len & 0x7B & X-type & Byte copy loop (microcoded) \\
CPYW dst, src, len & 0x7C & X-type & Word copy loop (microcoded) \\
MEMFILL dst, val, len & 0x7D & X-type & Memory fill loop (microcoded) \\ \bottomrule
\end{longtable}

\subsection{System \& Privileged (15 instructions)}

\begin{longtable}{@{}L{2.8cm}L{1.5cm}L{1.7cm}L{4.5cm}@{}}
\toprule
\textbf{Instruction} & \textbf{Opcode} & \textbf{Encoding} & \textbf{Micro-Operations} \\ \midrule
\endfirsthead
\toprule
\textbf{Instruction} & \textbf{Opcode} & \textbf{Encoding} & \textbf{Micro-Operations} \\ \midrule
\endhead
HLT & 0x7E & S-type & uHALT (stop execution) \\
RESET & 0x7F & S-type & uRESET (soft reset) \\
RDCR Rdst, CRn & 0x80 & S-type & uREAD\_CTRL\_REG Rdst,CRn \\
WRCR CRn, Rsrc & 0x81 & S-type & uWRITE\_CTRL\_REG CRn,Rsrc \\
SYSCALL & 0x82 & S-type & uSYSCALL (trap to kernel) \\
SYSRET & 0x83 & S-type & uSYSRET (return from kernel) \\
IRET & 0x84 & S-type & uIRET (return from interrupt) \\
CLI & 0x85 & S-type & uCLEAR\_IF (disable interrupts) \\
STI & 0x86 & S-type & uSET\_IF (enable interrupts) \\
WFI & 0x87 & S-type & uWAIT\_FOR\_INTERRUPT \\
RDTIME Rdst & 0x88 & S-type & uREAD\_TIMER Rdst \\
RDTS Rdst & 0x89 & S-type & uREAD\_TIMESTAMP Rdst \\
SLEEP \#imm & 0x8A & S-type + imm & uSLEEP \#imm (power mode) \\
GETPID Rdst & 0x8B & S-type & uGET\_PROCESS\_ID Rdst \\
GETTID Rdst & 0x8C & S-type & uGET\_THREAD\_ID Rdst \\ \bottomrule
\end{longtable}

\subsection{Stack \& Call Frame (10 instructions)}

\begin{longtable}{@{}L{2.8cm}L{1.5cm}L{1.7cm}L{4.5cm}@{}}
\toprule
\textbf{Instruction} & \textbf{Opcode} & \textbf{Encoding} & \textbf{Micro-Operations} \\ \midrule
\endfirsthead
\toprule
\textbf{Instruction} & \textbf{Opcode} & \textbf{Encoding} & \textbf{Micro-Operations} \\ \midrule
\endhead
PUSH Rsrc & 0x8D & S-type & uDEC SP; uSTORE [SP],Rsrc \\
POP Rdst & 0x8E & S-type & uLOAD Rdst,[SP]; uINC SP \\
PUSHA & 0x8F & X-type & Push all GPRs (R0-R15, microcoded) \\
POPA & 0x90 & X-type & Pop all GPRs (R15-R0, microcoded) \\
ENTER \#framesize & 0x91 & X-type & uPUSH BP; uMOV BP,SP; uSUB SP,size \\
LEAVE & 0x92 & X-type & uMOV SP,BP; uPOP BP \\
PUSHF & 0x93 & S-type & uDEC SP; uSTORE [SP],FLAGS \\
POPF & 0x94 & S-type & uLOAD FLAGS,[SP]; uINC SP \\
PUSHQ regpair & 0x95 & S-type & Push register pair (64-bit) \\
POPQ regpair & 0x96 & S-type & Pop register pair (64-bit) \\ \bottomrule
\end{longtable}

\subsection{SIMD / Vector (15 instructions)}

\begin{longtable}{@{}L{2.8cm}L{1.5cm}L{1.7cm}L{4.5cm}@{}}
\toprule
\textbf{Instruction} & \textbf{Opcode} & \textbf{Encoding} & \textbf{Micro-Operations} \\ \midrule
\endfirsthead
\toprule
\textbf{Instruction} & \textbf{Opcode} & \textbf{Encoding} & \textbf{Micro-Operations} \\ \midrule
\endhead
VADD Vd, Va, Vb & 0x97 & V-type & Per-lane: uADD Vd[i],Va[i],Vb[i] \\
VSUB Vd, Va, Vb & 0x98 & V-type & Per-lane: uSUB Vd[i],Va[i],Vb[i] \\
VMUL Vd, Va, Vb & 0x99 & V-type & Per-lane: uMUL Vd[i],Va[i],Vb[i] \\
VDIV Vd, Va, Vb & 0x9A & V-type & Per-lane: uDIV Vd[i],Va[i],Vb[i] \\
VAND Vd, Va, Vb & 0x9B & V-type & Per-lane: uAND Vd[i],Va[i],Vb[i] \\
VOR Vd, Va, Vb & 0x9C & V-type & Per-lane: uOR Vd[i],Va[i],Vb[i] \\
VXOR Vd, Va, Vb & 0x9D & V-type & Per-lane: uXOR Vd[i],Va[i],Vb[i] \\
VSHL Vd, Va, \#imm & 0x9E & V-type + imm & Per-lane: uSHL Vd[i],Va[i],imm \\
VSHR Vd, Va, \#imm & 0x9F & V-type + imm & Per-lane: uSHR Vd[i],Va[i],imm \\
VDUP Vd, Va & 0xA0 & V-type & Broadcast: Vd[i]=Va[0] for all i \\
VLD Vd, [addr] & 0xA1 & V-type + M & Load vector from memory \\
VST [addr], Vd & 0xA2 & V-type + M & Store vector to memory \\
VABS Vd, Va & 0xA3 & V-type & Per-lane: uABS Vd[i],Va[i] \\
VMAX Vd, Va, Vb & 0xA4 & V-type & Per-lane: uMAX Vd[i],Va[i],Vb[i] \\
VMIN Vd, Va, Vb & 0xA5 & V-type & Per-lane: uMIN Vd[i],Va[i],Vb[i] \\ \bottomrule
\end{longtable}

\subsection{Floating Point (10 instructions)}

\begin{longtable}{@{}L{2.8cm}L{1.5cm}L{1.7cm}L{4.5cm}@{}}
\toprule
\textbf{Instruction} & \textbf{Opcode} & \textbf{Encoding} & \textbf{Micro-Operations} \\ \midrule
\endfirsthead
\toprule
\textbf{Instruction} & \textbf{Opcode} & \textbf{Encoding} & \textbf{Micro-Operations} \\ \midrule
\endhead
FADD Fd, Fa, Fb & 0xA6 & R-type (FP) & fADD Fd,Fa,Fb (IEEE 754) \\
FSUB Fd, Fa, Fb & 0xA7 & R-type (FP) & fSUB Fd,Fa,Fb (IEEE 754) \\
FMUL Fd, Fa, Fb & 0xA8 & R-type (FP) & fMUL Fd,Fa,Fb (IEEE 754) \\
FDIV Fd, Fa, Fb & 0xA9 & R-type (FP) & fDIV Fd,Fa,Fb (IEEE 754) \\
FSQRT Fd, Fa & 0xAA & R-type (FP) & fSQRT Fd,Fa (square root) \\
FCMP Fa, Fb & 0xAB & R-type (FP) & fCMP Fa,Fb (set FP flags) \\
FCVTI Fd, Ra & 0xAC & R-type (conv) & Convert integer Ra to float Fd \\
FCVTS Ra, Fd & 0xAD & R-type (conv) & Convert float Fd to integer Ra \\
FNEG Fd, Fa & 0xAE & R-type (FP) & fNEG Fd,Fa (negate) \\
FABS Fd, Fa & 0xAF & R-type (FP) & fABS Fd,Fa (absolute value) \\ \bottomrule
\end{longtable}

\section{Micro-Operations Architecture}

\subsection{Hybrid CISC/RISC Design Philosophy}

The SeaBird ISA employs a hybrid architecture combining a CISC frontend with a RISC backend, similar to modern processors from Intel and AMD. This design provides:

\begin{itemize}
    \item \textbf{CISC Frontend:} Expressive, compact instructions visible to programmers and compilers
    \item \textbf{RISC Backend:} Simple, uniform micro-operations (µops) executed by the CPU pipeline
    \item \textbf{Performance Benefits:} Combines code density of CISC with execution efficiency of RISC
\end{itemize}

\textbf{Execution Flow:}

\begin{enumerate}
    \item Programmer writes CISC-style instructions (e.g., LOAD R1, [R2+R3*4+32])
    \item Frontend decoder translates to RISC-like micro-ops
    \item Backend execution engine processes micro-ops through standard pipeline stages
    \item Micro-ops can be reordered, fused, and optimized independently
\end{enumerate}

This approach enables:
\begin{itemize}
    \item Out-of-order execution
    \item Micro-op fusion and optimization
    \item Superscalar dispatch
    \item Register renaming at the µop level
    \item Cache and branch prediction at µop granularity
\end{itemize}

\subsection{Micro-Op Decomposition Rules}

The following rules govern how CISC instructions decompose into RISC-like micro-operations:

\subsubsection{Rule 1: Single-Operation Instructions $\rightarrow$ 1 µop}

Simple register-to-register operations map directly to single micro-ops.

\begin{lstlisting}
; CISC Instruction -> Micro-Op
ADD R1, R2, R3     -> uADD R1, R2, R3
XOR R4, R5         -> uXOR R4, R5
NEG R3             -> uNEG R3
MOV R6, R7         -> uMOV R6, R7
\end{lstlisting}

\subsubsection{Rule 2: Memory Operations $\rightarrow$ 2--3 µops}

Load/store with simple addressing requires address calculation plus memory access.

\begin{lstlisting}
; CISC Instruction
LOAD R1, [R2+16]

; Micro-Op Breakdown
uADD  Rtmp, R2, #16       ; compute effective address
uLOAD R1, [Rtmp]          ; perform load
\end{lstlisting}

\begin{lstlisting}
; CISC Instruction
STORE [R3+8], R4

; Micro-Op Breakdown
uADD   Rtmp, R3, #8       ; compute effective address
uSTORE [Rtmp], R4         ; perform store
\end{lstlisting}

\subsubsection{Rule 3: Complex Addressing Modes $\rightarrow$ 3--5 µops}

Complex addressing with base + scaled index + displacement requires multiple steps.

\begin{lstlisting}
; CISC Instruction
LOAD R1, [R2+R3*4+32]

; Micro-Op Breakdown
uMUL  Rtmp1, R3, #4       ; scale index by 4
uADD  Rtmp2, R2, Rtmp1    ; add base to scaled index
uADD  Rtmp2, Rtmp2, #32   ; add displacement
uLOAD R1, [Rtmp2]         ; perform load
\end{lstlisting}

\subsubsection{Rule 4: Arithmetic with Immediate $\rightarrow$ 1--2 µops}

Immediate arithmetic typically maps to single µop if immediate fits in encoding.

\begin{lstlisting}
; CISC Instruction
ADDI R1, R2, #7

; Micro-Op Breakdown
uADD R1, R2, #7           ; single µop with embedded immediate
\end{lstlisting}

For large immediates:

\begin{lstlisting}
; CISC Instruction
ADDI R1, R2, #0xDEADBEEF

; Micro-Op Breakdown
uMOVI Rtmp, #0xDEADBEEF   ; load large immediate
uADD  R1, R2, Rtmp        ; perform addition
\end{lstlisting}

\subsubsection{Rule 5: CISC-Style Compound Operations $\rightarrow$ 2--6 µops}

Complex operations combining multiple actions decompose into sequential µops.

\begin{lstlisting}
; CISC Instruction
PUSH [R3+8]

; Micro-Op Breakdown
uADD   Rtmp1, R3, #8      ; compute source address
uLOAD  Rtmp2, [Rtmp1]     ; load value from memory
uADD   RSP, RSP, #-8      ; adjust stack pointer
uSTORE [RSP], Rtmp2       ; store to stack
\end{lstlisting}

\begin{lstlisting}
; CISC Instruction
CALL label

; Micro-Op Breakdown
uMOVI  Rtmp1, return_addr ; load return address
uSTORE [RSP-8], Rtmp1     ; push return address
uADD   RSP, RSP, #-8      ; adjust stack pointer
uJMP   label              ; jump to target
\end{lstlisting}

\subsubsection{Rule 6: Branches $\rightarrow$ 1 µop}

Branch instructions map directly to conditional or unconditional µop branches.

\begin{lstlisting}
; CISC Instruction -> Micro-Op
JZ  label          -> uBRANCH_IF_ZERO PC+offset
JMP target         -> uJMP PC+offset
RET                -> uRET (implicit load+jump)
\end{lstlisting}

\subsubsection{Rule 7: Vector Operations $\rightarrow$ N µops}

Vector operations decompose into per-lane µops based on vector width.

\begin{lstlisting}
; CISC Instruction
VADD V0, V1, V2

; Micro-Op Breakdown (for 4-lane vector)
uADD V0.lane0, V1.lane0, V2.lane0
uADD V0.lane1, V1.lane1, V2.lane1
uADD V0.lane2, V1.lane2, V2.lane2
uADD V0.lane3, V1.lane3, V2.lane3
\end{lstlisting}

\textbf{Note:} Vector µops may execute in parallel on SIMD execution units.

\subsubsection{Rule 8: Microcoded Instructions $\rightarrow$ Variable µops}

Special or rarely-used instructions may invoke microcode sequences.

\begin{lstlisting}
; CISC Instruction
PUSHA                     ; push all GPRs

; Micro-Op Breakdown (microcoded)
uSTORE [RSP-8],  R0
uSTORE [RSP-16], R1
uSTORE [RSP-24], R2
...
uSTORE [RSP-248], R31
uADD RSP, RSP, #-256
\end{lstlisting}

\subsection{Micro-Op Execution Properties}

\subsubsection{Micro-Op Characteristics}

All micro-ops share common properties enabling efficient backend execution:

\begin{itemize}
    \item \textbf{Fixed-format:} Uniform µop structure simplifies decode and dispatch
    \item \textbf{Single-operation:} Each µop performs one elementary action
    \item \textbf{Explicit operands:} No implicit side effects (except flags)
    \item \textbf{Register-based:} Temporary registers (T0--T3) used for intermediate values
    \item \textbf{Predictable latency:} Most µops have 1-cycle execution time
\end{itemize}

\subsubsection{Backend Optimization Opportunities}

The µop stream enables aggressive backend optimizations:

\begin{itemize}
    \item \textbf{Micro-op fusion:} Adjacent µops combined into single operation
    \item \textbf{Register renaming:} Eliminate false dependencies
    \item \textbf{Out-of-order execution:} Independent µops execute in parallel
    \item \textbf{Speculative execution:} Branch prediction at µop granularity
    \item \textbf{Memory disambiguation:} Reorder loads/stores when safe
\end{itemize}

\subsubsection{Example: Full Instruction Pipeline}

\begin{lstlisting}
; Programmer writes:
ADD R1, R2, [R3+R4*8+64]

; Frontend decodes to micro-ops:
uMUL  T0, R4, #8          ; cycle 1: scale index
uADD  T1, R3, T0          ; cycle 2: add base
uADD  T1, T1, #64         ; cycle 3: add displacement
uLOAD T2, [T1]            ; cycle 4-6: memory access
uADD  R1, R2, T2          ; cycle 7: final addition

; Backend may optimize to:
; - Fuse address calculation µops
; - Execute uMUL and other ops out-of-order
; - Rename T0, T1, T2 to physical registers
; - Overlap memory access with other work
\end{lstlisting}

\section{Instruction Set Reference}

The SeaBird ISA defines 160 instructions across 10 functional categories. Each instruction may decompose into one or more micro-ops during execution.

\subsection{Data Movement (24 instructions)}

\subsubsection{Basic Moves}

\begin{lstlisting}
MOV     dst, src          ; register or memory move
MOVI    dst, imm          ; move immediate
MOVZX   dst, src          ; zero-extend
MOVSX   dst, src          ; sign-extend
MOVHI   dst, src          ; move high half
MOVLO   dst, src          ; move low half
MOVSWP  dst, src          ; byte-swap endian
\end{lstlisting}

\subsubsection{Memory Access}

\begin{lstlisting}
LD      dst, [addr]       ; load
ST      [addr], src       ; store
LDI     dst, imm(addr)    ; load from immediate address
LDB     dst, [addr]       ; load byte
LDH     dst, [addr]       ; load halfword
LDW     dst, [addr]       ; load word
LDQ     dst, [addr]       ; load quadword
STB     [addr], src       ; store byte
STH     [addr], src       ; store halfword
STW     [addr], src       ; store word
STQ     [addr], src       ; store quadword
\end{lstlisting}

\subsubsection{Special Variants}

\begin{lstlisting}
LDP     dst1, dst2, [addr]      ; load pair
STP     [addr], src1, src2      ; store pair
LEA     dst, [addr]             ; compute effective address
LEAS    dst, [base + index*scale + imm]
XCHG    dst, src                ; atomic swap
\end{lstlisting}

\subsection{Arithmetic (26 instructions)}

\subsubsection{Integer Arithmetic}

\begin{lstlisting}
ADD     r, v              ; add
ADDI    r, imm            ; add immediate
SUB     r, v              ; subtract
SUBI    r, imm            ; subtract immediate
MUL     dst, src          ; multiply
MULI    dst, imm          ; multiply immediate
DIV     dst, src          ; divide
DIVI    dst, imm          ; divide immediate
MOD     dst, src          ; modulo
MODI    dst, imm          ; modulo immediate
NEG     r                 ; negate
INC     r                 ; increment
DEC     r                 ; decrement
\end{lstlisting}

\subsubsection{Extended / 64-bit}

\begin{lstlisting}
MULH    dst, src          ; high-half multiply
UMUL    dst, src          ; unsigned multiply
UDIV    dst, src          ; unsigned divide
\end{lstlisting}

\subsubsection{Saturating Arithmetic}

\begin{lstlisting}
ADDS    dst, src          ; signed saturating add
ADDU    dst, src          ; unsigned saturating add
SUBS    dst, src          ; signed saturating subtract
SUBU    dst, src          ; unsigned saturating subtract
\end{lstlisting}

\subsubsection{Bit Arithmetic}

\begin{lstlisting}
ABS     dst, src          ; absolute value
CLZ     dst, src          ; count leading zeros
CTZ     dst, src          ; count trailing zeros
POPC    dst, src          ; popcount
\end{lstlisting}

\subsection{Logic \& Bit Operations (18 instructions)}

\begin{lstlisting}
AND     r, v              ; logical AND
OR      r, v              ; logical OR
XOR     r, v              ; logical XOR
NOT     r                 ; logical NOT
NAND    r, v              ; NAND
NOR     r, v              ; NOR
XNOR    r, v              ; XNOR
SHL     r, v              ; shift left
SHR     r, v              ; shift right
SAR     r, v              ; shift arithmetic right
ROL     r, v              ; rotate left
ROR     r, v              ; rotate right
BSET    r, bit            ; bit set
BCLR    r, bit            ; bit clear
BTOG    r, bit            ; bit toggle
BTST    dst, bit          ; test bit
MASK    dst, src, mask    ; apply mask
EXT     dst, src, bitrange ; extract bit range
\end{lstlisting}

\subsection{Comparison (10 instructions)}

\begin{lstlisting}
CMP     a, b              ; compare (sets flags)
CMPI    a, imm            ; compare immediate
CMPS    a, b              ; signed compare
CMPU    a, b              ; unsigned compare
TST     a, b              ; logical test
TSTI    a, imm            ; test immediate
MAX     dst, src          ; maximum
MIN     dst, src          ; minimum
SLT     dst, src          ; set if less than
SGT     dst, src          ; set if greater than
\end{lstlisting}

\subsection{Branching \& Control Flow (20 instructions)}

\subsubsection{Unconditional}

\begin{lstlisting}
JMP     lbl               ; jump to label
JMPA    addr              ; jump to absolute address
CALL    lbl               ; call subroutine
CALLA   addr              ; call absolute address
RET                       ; return from subroutine
\end{lstlisting}

\subsubsection{Conditional (flag-based)}

\begin{lstlisting}
JE      lbl               ; jump if equal
JNE     lbl               ; jump if not equal
JG      lbl               ; jump if greater
JGE     lbl               ; jump if greater or equal
JL      lbl               ; jump if less
JLE     lbl               ; jump if less or equal
JC      lbl               ; jump if carry
JNC     lbl               ; jump if not carry
JO      lbl               ; jump if overflow
JNO     lbl               ; jump if not overflow
JS      lbl               ; jump if sign
JNS     lbl               ; jump if not sign
\end{lstlisting}

\subsubsection{Register-based}

\begin{lstlisting}
JZR     reg, lbl          ; jump if register zero
JNZR    reg, lbl          ; jump if register not zero
BRR     reg               ; jump via register
\end{lstlisting}

\subsubsection{Special}

\begin{lstlisting}
TRAP    imm               ; software trap
YIELD                     ; yield execution
\end{lstlisting}

\subsection{Memory / Addressing Extensions (12 instructions)}

\begin{lstlisting}
PREFETCH    [addr]        ; prefetch cache line
FLUSH       [addr]        ; flush cache line
INVIC       [addr]        ; invalidate icache line
INVDC       [addr]        ; invalidate dcache line
LDX         dst, [addr+]  ; load and post-increment
STX         [addr+], src  ; store and post-increment
LDN         dst, [addr-]  ; load and post-decrement
STN         [addr-], src  ; store and post-decrement
MOVM        dst, mask     ; masked move
CPYB        dst, src, len ; copy bytes
CPYW        dst, src, len ; copy words
MEMFILL     dst, value, len ; fill memory
\end{lstlisting}

\subsection{System \& Privileged (15 instructions)}

\begin{lstlisting}
HLT                       ; halt processor
RESET                     ; system reset
RDCR    dst, crN          ; read control register
WRCR    crN, src          ; write control register
SYSCALL                   ; system call
SYSRET                    ; return from syscall
IRET                      ; interrupt return
CLI                       ; clear interrupts
STI                       ; set interrupts
WFI                       ; wait for interrupt
RDTIME  dst               ; read time
RDTS    dst               ; read timestamp
SLEEP   imm               ; sleep for duration
GETPID                    ; get process ID
GETTID                    ; get thread ID
\end{lstlisting}

\subsection{Stack \& Call Frame (10 instructions)}

\begin{lstlisting}
PUSH    r                 ; push register
POP     r                 ; pop register
PUSHA                     ; push all GPRs
POPA                      ; pop all GPRs
ENTER   imm               ; create stack frame
LEAVE                     ; destroy stack frame
PUSHF                     ; push flags
POPF                      ; pop flags
PUSHQ   regpair           ; push register pair
POPQ    regpair           ; pop register pair
\end{lstlisting}

\subsection{SIMD / Vector (15 instructions)}

Assuming 128-bit or 256-bit V registers:

\begin{lstlisting}
VADD    vd, va, vb        ; vector add
VSUB    vd, va, vb        ; vector subtract
VMUL    vd, va, vb        ; vector multiply
VDIV    vd, va, vb        ; vector divide
VAND    vd, va, vb        ; vector AND
VOR     vd, va, vb        ; vector OR
VXOR    vd, va, vb        ; vector XOR
VSHL    vd, va, imm       ; vector shift left
VSHR    vd, va, imm       ; vector shift right
VDUP    vd, va            ; vector duplicate
VLD     vd, [addr]        ; vector load
VST     [addr], vd        ; vector store
VABS    vd, va            ; vector absolute value
VMAX    vd, va, vb        ; vector maximum
VMIN    vd, va, vb        ; vector minimum
\end{lstlisting}

\subsection{Floating Point (10 instructions)}

\begin{lstlisting}
FADD    fd, fa, fb        ; floating-point add
FSUB    fd, fa, fb        ; floating-point subtract
FMUL    fd, fa, fb        ; floating-point multiply
FDIV    fd, fa, fb        ; floating-point divide
FSQRT   fd, fa            ; floating-point square root
FCMP    fa, fb            ; floating-point compare
FCVTI   fd, ra            ; convert int to float
FCVTS   ra, fd            ; convert float to int
FNEG    fd                ; floating-point negate
FABS    fd                ; floating-point absolute value
\end{lstlisting}

\textbf{Total: 160 instructions}

\section{Assembler Directives}

This section documents all assembler directives supported by the toolchain. Directives control code organization, symbol visibility, data layout, optimization behavior, and architecture-specific features.

\subsection{Code \& Data Layout Directives}

\begin{description}
\item[\texttt{.text}] Begin a code section. Assembler switches to instruction-emission mode.

\item[\texttt{.data}] Begin a data section. Used for static variables, lookup tables, and buffers.

\item[\texttt{.bss}] Begin an uninitialized data section. Reserves space without emitting bytes.

\item[\texttt{.align \textit{n}}] Align the next emitted object to an \textit{n}-byte boundary. Particularly beneficial for aligning instructions or micro-tables in CISC frontend code.

\item[\texttt{.org \textit{addr}}] Set the absolute location counter to \textit{addr}. Primarily used for firmware, vector tables, and boot ROM.
\end{description}

\subsection{Symbol \& Label Directives}

\begin{description}
\item[\texttt{.global \textit{symbol}}] Export a symbol for use by other modules.

\item[\texttt{.extern \textit{symbol}}] Declare a symbol to be resolved by the linker.

\item[\texttt{.equ \textit{name}, \textit{value}}] Create a constant. Equivalent to \texttt{\#define} in C.

\item[\texttt{.set \textit{name}, \textit{value}}] Same as \texttt{.equ}, but may be reassigned later.

\item[\texttt{.macro \textit{name} \textit{args...}}] Begin a macro definition.

\item[\texttt{.endm}] End a macro definition.
\end{description}

\subsection{Data Definition Directives}

\begin{description}
\item[\texttt{.byte \textit{b1}, \textit{b2}, ...}] Emit raw 8-bit bytes.

\item[\texttt{.word \textit{w1}, \textit{w2}, ...}] Emit 16-bit values.

\item[\texttt{.dword \textit{d1}, \textit{d2}, ...}] Emit 32-bit values.

\item[\texttt{.qword \textit{q1}, \textit{q2}, ...}] Emit 64-bit values.

\item[\texttt{.float \textit{f1}, \textit{f2}, ...}] Emit 32-bit IEEE floating-point values.

\item[\texttt{.double \textit{d1}, \textit{d2}, ...}] Emit 64-bit IEEE floating-point values.

\item[\texttt{.ascii "\textit{text}"}] Emit raw ASCII bytes without null terminator.

\item[\texttt{.asciz "\textit{text}"} / \texttt{.string "\textit{text}"}] Emit a null-terminated string.

\item[\texttt{.space \textit{n}}] Reserve \textit{n} bytes of space (zero-initialized or uninitialized).
\end{description}

\subsection{Control Directives}

\begin{description}
\item[\texttt{.include "\textit{file}"}] Insert another source file at this location.

\item[\texttt{.ifdef \textit{SYM}} / \texttt{.ifndef \textit{SYM}}] Begin conditional assembly block.

\item[\texttt{.endif}] End conditional assembly block.
\end{description}

\subsection{Architecture, Mode \& Optimization Control}

These directives are critical for controlling behavior in the hybrid ISA environment.

\begin{description}
\item[\texttt{.cpu \textit{model}}] Declare target CPU model or extension level. Enables assembler warnings if instructions are not supported on the specified model.

\item[\texttt{.arch \textit{ext1}, \textit{ext2}, ...}] Enable optional ISA extensions (e.g., \texttt{VECTOR}, \texttt{SYS}, \texttt{CRYPTO}).

\item[\texttt{.mopt \textit{level}}] Set assembler micro-op optimization level:
\begin{itemize}
\item \texttt{0} -- No micro-fusion (debug-friendly)
\item \texttt{1} -- Allow simple micro-op fusion
\item \texttt{2} -- Full aggressive micro-crunching mode
\end{itemize}

\item[\texttt{.mode \textit{MODE\_NAME}}] Switch code-generation mode. Supported modes include:
\begin{itemize}
\item \texttt{CLOWNFISH} -- 16-bit legacy mode
\item \texttt{TETRA} -- 32-bit base mode
\item \texttt{DRAGONET} -- 64-bit Meisei mode
\item \texttt{DROPLET} -- 128-bit hyper mode
\item \texttt{REALTIME} -- Disallow slow CISC expansions; enforce micro-op minimality
\end{itemize}

\item[\texttt{.uops on/off}] Show generated micro-ops in listing output.
\end{description}

\subsection{Debugging, Listing \& Metadata}

\begin{description}
\item[\texttt{.file "\textit{name}"}] Set debug filename for debugging symbols.

\item[\texttt{.line \textit{n}}] Force encoding of source line number \textit{n} for debugging.

\item[\texttt{.section \textit{name}}] Create or switch to user-defined sections. Useful for custom loaders or non-standard binary formats.

\item[\texttt{.message "\textit{text}"}] Print informational message during assembly.

\item[\texttt{.warning "\textit{text}"}] Print warning message (assembly continues).

\item[\texttt{.error "\textit{text}"}] Emit fatal error and halt assembly.
\end{description}

\subsection{Special Optimization Directives}

These directives provide fine-grained control over assembler optimization behavior.

\begin{description}
\item[\texttt{.hint microburst}] Mark the following block as performance-critical. Instructs the optimizer to apply aggressive optimizations.

\item[\texttt{.prefetch \textit{addr}}] Emit an embedded prefetch hint micro-op for address \textit{addr}.

\item[\texttt{.autovec on/off}] Enable or disable automatic loop vectorization by the assembler.

\item[\texttt{.trustme}] Disable \textbf{all} safety checks. Assembler will accept code without validation. \textit{Use with extreme caution.}

\item[\texttt{.sleepdeprived}] Convenience directive equivalent to:
\begin{verbatim}
    .mopt 2
    .uops on
    .autovec on
\end{verbatim}
Enables maximum optimization and diagnostic output. Named for late-night coding sessions.
\end{description}

\section{Example Instructions With Micro-Op Breakdown}

\begin{longtable}{@{}L{2.8cm}L{7cm}@{}}
\toprule
\textbf{Instruction} & \textbf{Micro-Ops} \\ \midrule
\endfirsthead
\toprule
\textbf{Instruction} & \textbf{Micro-Ops} \\ \midrule
\endhead
ADD R1,R2,R3 & uADD R1,R2,R3 \\
LOAD R1,[R2+R3*4+32] & uMUL Rtmp,R3,\#4; uADD Rtmp,R2,Rtmp; uADD Rtmp,Rtmp,\#32; uLOAD R1,[Rtmp] \\
CALL label & uMOVI TMP1,return\_addr; uSTORE [SP-8],TMP1; uALU SP,SP,-8,ADD; uJMP label \\
RET & uLOAD TMP1,[SP]; uALU SP,SP,8,ADD; uJMP TMP1 \\
VADD V0,V1,V2 & loop over lanes: uADD lane0,lane1,lane2 \\ \bottomrule
\end{longtable}

\section{Instruction Latency \& Performance Characteristics}

\subsection{Typical Instruction Latencies}

Performance characteristics for a reference implementation with out-of-order execution:

\begin{table}[h]
\centering
\begin{tabular}{@{}lll@{}}
\toprule
\textbf{Operation Type} & \textbf{Latency} & \textbf{Throughput} \\ \midrule
Single µop (ADD, MOV, XOR) & 1 cycle & 4/cycle \\
Memory load (L1 hit) & 4 cycles & 2/cycle \\
Memory store (L1 hit) & 1 cycle & 2/cycle \\
Integer multiply & 3 cycles & 1/cycle \\
Integer divide & 10--20 cycles & 1/10 cycles \\
Floating-point add/sub & 3 cycles & 2/cycle \\
Floating-point multiply & 4 cycles & 2/cycle \\
Floating-point divide & 10--15 cycles & 1/10 cycles \\
Vector ops (per lane) & 1--4 cycles & varies \\
Microcoded instructions & 3--20 cycles & varies \\
Branch (predicted) & 1 cycle & 2/cycle \\
Branch (mispredicted) & 15--20 cycles & penalty \\ \bottomrule
\end{tabular}
\end{table}

\subsection{Cache and Memory Hierarchy}

Typical implementation memory subsystem:

\begin{table}[h]
\centering
\begin{tabular}{@{}llll@{}}
\toprule
\textbf{Level} & \textbf{Size} & \textbf{Latency} & \textbf{Bandwidth} \\ \midrule
L1 Data Cache & 32--64 KB & 4 cycles & 64 GB/s \\
L1 Instruction Cache & 32--64 KB & 4 cycles & 64 GB/s \\
L2 Unified Cache & 256--512 KB & 12 cycles & 32 GB/s \\
L3 Unified Cache & 4--16 MB & 40 cycles & 16 GB/s \\
Main Memory (DDR4) & 8--64 GB & 100--300 cycles & 25 GB/s \\ \bottomrule
\end{tabular}
\end{table}

\subsection{Pipeline Characteristics}

\begin{itemize}
    \item \textbf{Pipeline depth:} 14--20 stages (typical superscalar implementation)
    \item \textbf{Fetch width:} 4--6 instructions per cycle
    \item \textbf{Decode width:} 4--6 instructions → 8--12 µops per cycle
    \item \textbf{Execution width:} 8--10 µops per cycle
    \item \textbf{Retire width:} 4--6 instructions per cycle
    \item \textbf{ROB size:} 192--256 µops (reorder buffer)
    \item \textbf{Branch prediction:} Two-level adaptive with 4K--8K entry BTB
\end{itemize}

\section{FLAGS Register}

\subsection{Status Flags}

The FLAGS register contains condition codes and processor state. Bits are arranged as follows:

\begin{longtable}{@{}llll@{}}
\toprule
\textbf{Bit} & \textbf{Name} & \textbf{Symbol} & \textbf{Description} \\ \midrule
\endfirsthead
\toprule
\textbf{Bit} & \textbf{Name} & \textbf{Symbol} & \textbf{Description} \\ \midrule
\endhead
0 & Carry & CF & Carry out of MSB or borrow \\
1 & Reserved & -- & Must be 0 \\
2 & Parity & PF & Even parity of low byte \\
3 & Reserved & -- & Must be 0 \\
4 & Auxiliary Carry & AF & Carry from bit 3 to bit 4 \\
5 & Reserved & -- & Must be 0 \\
6 & Zero & ZF & Result is zero \\
7 & Sign & SF & Sign bit of result (MSB) \\
8 & Trap & TF & Single-step debugging \\
9 & Interrupt Enable & IF & Interrupts enabled \\
10 & Direction & DF & String operation direction \\
11 & Overflow & OF & Signed arithmetic overflow \\
12--13 & I/O Privilege Level & IOPL & I/O privilege (0--3) \\
14 & Nested Task & NT & Nested task flag \\
15 & Reserved & -- & Must be 0 \\
16 & Resume & RF & Resume flag \\
17 & Virtual 8086 & VM & Virtual 8086 mode \\
18 & Alignment Check & AC & Alignment check enabled \\
19 & Virtual Interrupt & VIF & Virtual interrupt flag \\
20 & Virtual Int Pending & VIP & Virtual interrupt pending \\
21 & ID & ID & CPUID support \\
22--31 & Reserved & -- & Must be 0 \\ \bottomrule
\end{longtable}

\subsection{Flag Modification}

Instructions affect flags as follows:

\begin{itemize}
    \item \textbf{Arithmetic:} ADD, SUB, MUL, DIV modify CF, ZF, SF, OF, PF, AF
    \item \textbf{Logic:} AND, OR, XOR clear CF and OF, modify ZF, SF, PF
    \item \textbf{Shift/Rotate:} Modify CF and OF
    \item \textbf{Compare:} CMP, TEST modify all arithmetic flags
    \item \textbf{Branches:} Read flags but do not modify
\end{itemize}

\section{Exception \& Interrupt Handling}

\subsection{Exception Vector Table}

The exception vector table begins at address 0xCAFE and contains handlers for system events:

\begin{longtable}{@{}llll@{}}
\toprule
\textbf{Vector} & \textbf{Mnemonic} & \textbf{Type} & \textbf{Description} \\ \midrule
\endfirsthead
\toprule
\textbf{Vector} & \textbf{Mnemonic} & \textbf{Type} & \textbf{Description} \\ \midrule
\endhead
0x00 & RESET & System & Power-on/reset \\
0x01 & DIV\_ZERO & Fault & Division by zero \\
0x02 & DEBUG & Trap & Debug/single-step \\
0x03 & NMI & Interrupt & Non-maskable interrupt \\
0x04 & BREAKPOINT & Trap & Software breakpoint \\
0x05 & OVERFLOW & Trap & Arithmetic overflow \\
0x06 & BOUND & Fault & Bounds check failure \\
0x07 & INVALID\_OP & Fault & Invalid opcode \\
0x08 & NO\_FPU & Fault & FPU not available \\
0x09 & DOUBLE\_FAULT & Abort & Double fault \\
0x0A & INVALID\_TSS & Fault & Invalid task segment \\
0x0B & SEG\_NOT\_PRESENT & Fault & Segment not present \\
0x0C & STACK\_FAULT & Fault & Stack segment fault \\
0x0D & GPF & Fault & General protection fault \\
0x0E & PAGE\_FAULT & Fault & Page fault \\
0x0F & RESERVED & -- & Reserved \\
0x10 & FPU\_ERROR & Fault & Floating-point error \\
0x11 & ALIGN\_CHECK & Fault & Alignment check \\
0x12 & MACHINE\_CHECK & Abort & Machine check \\
0x13 & SIMD\_ERROR & Fault & SIMD floating-point \\
0x14--0x1F & RESERVED & -- & Reserved for future use \\
0x20--0xFF & USER\_INT & Interrupt & External/software interrupts \\ \bottomrule
\end{longtable}

\subsection{Exception Handling Mechanism}

When an exception occurs:

\begin{enumerate}
    \item Current instruction address saved to CFAR (CISC Frontend Address Register)
    \item Current FLAGS saved to stack
    \item Current IP saved to stack
    \item IF flag cleared (interrupts disabled)
    \item Control transferred to exception handler via vector table
    \item Handler executes, performs recovery or termination
    \item Handler returns via IRET instruction
\end{enumerate}

\subsection{Privilege Level Transitions}

Exceptions may cause privilege level changes:

\begin{itemize}
    \item \textbf{User → Kernel:} On system call, page fault, or protection violation
    \item \textbf{Kernel → Kernel:} On nested exceptions
    \item \textbf{Kernel → User:} On IRET or SYSRET
\end{itemize}

Stack switching occurs automatically when privilege level changes, using kernel stack pointer from special register.

\section{Control Registers}

\subsection{System Control Registers}

Control registers (CR0--CR15) manage system-level functionality:

\begin{longtable}{@{}lll@{}}
\toprule
\textbf{Register} & \textbf{Name} & \textbf{Purpose} \\ \midrule
\endfirsthead
\toprule
\textbf{Register} & \textbf{Name} & \textbf{Purpose} \\ \midrule
\endhead
CR0 & System Control & Protected mode, paging, FPU control \\
CR1 & Reserved & Reserved for future use \\
CR2 & Page Fault Linear Address & Faulting address on page fault \\
CR3 & Page Directory Base & Physical address of page directory \\
CR4 & Extended Control & Additional CPU features \\
CR5--CR7 & Reserved & Reserved for future use \\
CR8 & Task Priority & Interrupt priority level \\
CR9--CR15 & Implementation-specific & Vendor-specific features \\ \bottomrule
\end{longtable}

\subsection{CR0 Bit Definitions}

\begin{table}[h]
\centering
\begin{tabular}{@{}lll@{}}
\toprule
\textbf{Bit} & \textbf{Name} & \textbf{Description} \\ \midrule
0 & PE & Protected Mode Enable \\
1 & MP & Monitor Coprocessor \\
2 & EM & Emulation (FPU emulation) \\
3 & TS & Task Switched \\
4 & ET & Extension Type \\
5 & NE & Numeric Error \\
16 & WP & Write Protect \\
18 & AM & Alignment Mask \\
29 & NW & Not Write-through \\
30 & CD & Cache Disable \\
31 & PG & Paging Enable \\ \bottomrule
\end{tabular}
\end{table}

\section{Memory Management}

\subsection{Paging Mechanism}

The SeaBird ISA supports hierarchical paging in Tetra and Dragonet modes:

\textbf{Tetra Mode (32-bit):}
\begin{itemize}
    \item Two-level page tables
    \item Page size: 4 KB (standard) or 4 MB (large pages)
    \item Virtual address space: 4 GB
    \item Page directory: 1024 entries
    \item Page table: 1024 entries per directory
\end{itemize}

\textbf{Dragonet Mode (64-bit):}
\begin{itemize}
    \item Four-level page tables (PML4)
    \item Page size: 4 KB, 2 MB (huge), or 1 GB (gigantic)
    \item Virtual address space: 256 TB (48-bit addresses)
    \item Canonical address form required
\end{itemize}

\subsection{Page Table Entry Format}

\begin{table}[h]
\centering
\begin{tabular}{@{}lll@{}}
\toprule
\textbf{Bits} & \textbf{Field} & \textbf{Description} \\ \midrule
0 & P & Present (page is in memory) \\
1 & R/W & Read/Write (0=read-only, 1=writable) \\
2 & U/S & User/Supervisor (0=kernel, 1=user) \\
3 & PWT & Page Write-Through \\
4 & PCD & Page Cache Disable \\
5 & A & Accessed (set by CPU) \\
6 & D & Dirty (written to) \\
7 & PS & Page Size (0=4KB, 1=large) \\
8 & G & Global page \\
9--11 & AVL & Available for OS use \\
12--51 & PFN & Physical Frame Number \\
52--62 & AVL & Available for OS use \\
63 & XD & Execute Disable (NX bit) \\ \bottomrule
\end{tabular}
\end{table}

\subsection{Translation Lookaside Buffer (TLB)}

\begin{itemize}
    \item \textbf{Data TLB:} 64--128 entries for load/store translations
    \item \textbf{Instruction TLB:} 64--128 entries for instruction fetch
    \item \textbf{Unified TLB (L2):} 512--1024 entries
    \item \textbf{Huge page support:} Separate entries for 2MB/1GB pages
    \item \textbf{PCID support:} Process Context ID for tagged TLB entries
\end{itemize}

\section{Programming Examples}

\subsection{Example 1: Simple Loop}

\begin{lstlisting}
; Calculate sum of array elements
; R1 = array pointer, R2 = count, R3 = sum

        MOVI    R3, #0              ; sum = 0
loop:
        LD      R4, [R1]            ; load array[i]
        ADD     R3, R3, R4          ; sum += array[i]
        ADDI    R1, #4              ; i++
        SUBI    R2, #1              ; count--
        JNZ     loop                ; if (count != 0) goto loop
        
        ; Result in R3
\end{lstlisting}

\subsection{Example 2: Function Call with Stack Frame}

\begin{lstlisting}
; Function: int add(int a, int b)
; Arguments: R0 = a, R1 = b
; Returns: R0 = result

add_function:
        ENTER   #16                 ; create stack frame (16 bytes)
        PUSH    R6                  ; save callee-saved register
        PUSH    R7
        
        MOV     R6, R0              ; save parameter a
        MOV     R7, R1              ; save parameter b
        ADD     R0, R6, R7          ; result = a + b
        
        POP     R7                  ; restore registers
        POP     R6
        LEAVE                       ; destroy stack frame
        RET                         ; return to caller

; Caller:
main:
        MOVI    R0, #10             ; first argument
        MOVI    R1, #20             ; second argument
        CALL    add_function        ; call function
        ; result in R0 (= 30)
\end{lstlisting}

\subsection{Example 3: Vector Operations}

\begin{lstlisting}
; Vector addition: C = A + B (4 elements each)
; V0 = vector A, V1 = vector B, V2 = result C

        VLD     V0, [R1]            ; load vector A
        VLD     V1, [R2]            ; load vector B
        VADD    V2, V0, V1          ; C = A + B
        VST     [R3], V2            ; store result
\end{lstlisting}

\subsection{Example 4: System Call}

\begin{lstlisting}
; System call: write(fd, buffer, count)
; Syscall number in R0, args in R1, R2, R3

        MOVI    R0, #1              ; syscall: write
        MOVI    R1, #1              ; fd = stdout
        LEA     R2, [message]       ; buffer pointer
        MOVI    R3, #13             ; count = 13 bytes
        SYSCALL                     ; invoke kernel
        ; return value in R0

message:
        .ascii  "Hello, World!"
\end{lstlisting}

\section{Implementation Notes}

\subsection{Design Considerations for Implementers}

\subsubsection{Frontend Decoder}

\begin{itemize}
    \item Complex instructions should decode to µop sequences in 1--2 cycles
    \item Maintain µop cache for frequently-decoded instruction sequences
    \item Support macro-op fusion (e.g., CMP + JZ → single µop)
\end{itemize}

\subsubsection{Execution Backend}

\begin{itemize}
    \item Minimum 4-wide superscalar for competitive performance
    \item Out-of-order execution with register renaming
    \item Unified reservation station or distributed scheduler
    \item Memory disambiguation for load/store reordering
\end{itemize}

\subsubsection{Branch Prediction}

\begin{itemize}
    \item Two-level adaptive predictor recommended
    \item BTB (Branch Target Buffer) with 4K--8K entries
    \item Return stack buffer with 16--32 entries
    \item Indirect branch predictor for virtual calls
\end{itemize}

\subsection{Compliance Levels}

Implementations may target different compliance levels:

\begin{itemize}
    \item \textbf{Level 0 (Minimal):} Clownfish mode only, single-issue, in-order
    \item \textbf{Level 1 (Basic):} Tetra mode, 2-way superscalar, simple branch prediction
    \item \textbf{Level 2 (Standard):} Dragonet mode, 4-way superscalar, out-of-order
    \item \textbf{Level 3 (High-Performance):} Full vector support, 6+ way, aggressive µop fusion
\end{itemize}

\section{Revision History}

\begin{table}[h]
\centering
\begin{tabular}{@{}llll@{}}
\toprule
\textbf{Version} & \textbf{Date} & \textbf{Author} & \textbf{Changes} \\ \midrule
1.0 & 2025-12-07 & SeaBird Team & Initial release \\ 1.1 & 2026-01-28 & Miyamii & Updated Directives\\ \bottomrule
\end{tabular}
\end{table}

\section{Acknowledgments}

The SeaBird ISA draws inspiration from several influential architectures:

\begin{itemize}
    \item x86-64 (Intel/AMD) - Binary encoding format, ModR/M/SIB structure
    \item ARM - Register naming conventions, clean instruction orthogonality
    \item RISC-V - Modern ISA design principles, extensibility
    \item IBM POWER - Hybrid CISC/RISC execution model
\end{itemize}

\section{Glossary}

\begin{description}
    \item[CISC] Complex Instruction Set Computer - Architecture with variable-length, multi-operation instructions
    \item[RISC] Reduced Instruction Set Computer - Architecture with fixed-length, single-operation instructions
    \item[µop (Micro-op)] Micro-operation - Internal RISC-like operation into which CISC instructions decompose
    \item[GPR] General-Purpose Register - Registers available for general computation (R0--R31)
    \item[MMU] Memory Management Unit - Hardware managing virtual-to-physical address translation
    \item[TLB] Translation Lookaside Buffer - Cache for page table entries
    \item[ROB] Reorder Buffer - Structure for out-of-order execution commit
    \item[BTB] Branch Target Buffer - Cache for branch target addresses
    \item[SIMD] Single Instruction Multiple Data - Parallel processing of vector data
    \item[ModR/M] Mode-Register-Memory byte - Encoding byte specifying operands
    \item[SIB] Scale-Index-Base byte - Encoding byte for complex addressing modes
\end{description}

\end{document}
